\documentclass[11pt]{article}

\usepackage[a4paper,margin=25mm,headheight=15pt]{geometry}
\usepackage[T1]{fontenc}
\usepackage[utf8]{inputenc}
\usepackage{lmodern}
\usepackage{microtype}
\usepackage{amsmath}
\usepackage{graphicx}
\usepackage{caption}
\usepackage{pdflscape}
\usepackage{xcolor}
\usepackage{fancyhdr}
\usepackage[hidelinks]{hyperref}

\definecolor{Accent}{HTML}{315B7D}
\definecolor{Muted}{HTML}{5B6570}

\hypersetup{pdftitle={Supplementary Figures: threshold sensitivity of the ancestry-sorting analyses},
  pdfauthor={Poikela et al.}}

\pagestyle{fancy}\fancyhf{}
\lhead{\small\color{Muted}Supplementary Figures}
\rhead{\small\color{Muted}Poikela et al.}
\cfoot{\small\thepage}
\renewcommand{\headrulewidth}{0.4pt}
\setlength{\parindent}{0pt}\setlength{\parskip}{0.6em}

\renewcommand{\thefigure}{S\arabic{figure}}
\newcommand{\Faq}{\textit{F.~aquilonia}}
\newcommand{\Fpo}{\textit{F.~polyctena}}

\title{\vspace{-3em}\textbf{\color{Accent}Supplementary Figures}\\[0.3em]
  \Large Threshold sensitivity of the ancestry-sorting analyses}
\author{Poikela et al.}
\date{}

\begin{document}
\maketitle
\vspace{-1.5em}

% ---- Methods / rationale -----------------------------------------------
The per-SNP sorting classification depends on two thresholds (Materials and
Methods): the replicate-population sorting threshold \(\tau\), the minimum
fraction of populations in which the dominant (unidirectional or bidirectional)
component must occur for a unit to be classified as sorted; and the
per-population near-fixation floor \(\phi\), the frequency the fixed-for
(major) allele must reach for a population to count as near-fixed. We evaluated
the sensitivity of the classification over a grid of
\(\tau \in \{0.5, 0.6, 0.7, 0.8\}\) and
\(\phi \in \{0.80, 0.85, 0.90, 0.95\}\). The upstream ancestry-diagnostic gate
(pooled-parental minor-allele frequency \(\ge 0.15\)) was held fixed, so the
denominator---the set of differentiated SNPs---is identical across every cell
of the grid and the resulting percentages are directly comparable. All other
conventions match the main analysis. The primary value \(\tau=0.5\) is the
logical minimum at which a sorting direction is shared by at least half of the
populations.

% ---- Figure S1: sorting-classification sweep ---------------------------
\begin{figure}[!ht]\centering
\includegraphics[width=\textwidth,trim=0 0 0 0.60in,clip]{../Figures/moduleA_sorting_sweep.png}
\caption{\textbf{Threshold sensitivity of the per-SNP sorting classification.}
Percentage of differentiated SNPs assigned to each sorting outcome as a function
of the sorting threshold \(\tau\) (\(x\)-axis) and the major-allele near-fixation
floor \(\phi\) (line colour). The differentiated SNP set and its parental
minor-allele-frequency gate were fixed across the grid. Panels show all sorted
SNPs and the unidirectional \Faq{}, unidirectional \Fpo{}, and bidirectional
classes. Increasing either threshold reduced the estimated prevalence of
sorting monotonically, from \(38.6\%\) under the least stringent combination to
\(1.3\%\) under the most stringent combination. At the primary operating point
(\(\tau = 0.5\), \(\phi = 0.90\)), \(25.4\%\) were classified as sorted.
Unidirectional sorting in either ancestry direction was consistently much more
common than bidirectional sorting, which approached zero as stringency
increased.}
\label{figS:sortingsweep}
\end{figure}

% ---- Figure S2: DI-dependent direction sweep ---------------------------
\clearpage
\begin{center}
\includegraphics[width=\textwidth,trim=0 0 0 0.60in,clip]{../Figures/moduleA_panelB_sweep.png}
\captionof{figure}{\textbf{Threshold sensitivity of the DI-dependent sorting direction.}
Fraction of unidirectional loci fixing toward \Faq{} across diagnostic-index
deciles (\(x\)-axis; low to high), for each combination of the sorting threshold
\(\tau\) (line colour) and the major-allele near-fixation floor \(\phi\)
(panels). The dashed horizontal line at \(0.5\) marks no directional bias:
points above it indicate a net lean toward \Faq{}, points below toward \Fpo{}.
DI decile boundaries were defined once from the threshold-independent set of
differentiated SNPs and held fixed across the grid. The loci contributing to
each fraction nevertheless vary because changing \(\tau\) or \(\phi\) changes
which loci meet the unidirectional-sorting criteria. Despite this change in the
qualifying set, the low-DI-toward-\Fpo{} / high-DI-toward-\Faq{} reversal seen
in main-text Figure~1B is preserved across the swept thresholds. The primary
operating point (\(\tau = 0.5\), \(\phi = 0.90\)) reproduces the main-text
curve exactly.}
\label{figS:panelBsweep}
\end{center}

% ---- Figure S3: direction-model coefficients vs sorting threshold ------
A complementary, architectural question is whether recombination predicts the
\emph{direction} of sorting beyond the diagnostic index. Among unidirectional
LD-reduced units we fit a logistic model,
P(\Faq)\,\(\sim\)\,diagnostic index \(+\) recombination \(+\) parental MAF
\(+\) \(\log\) block size, with all predictors standardised, and asked how its
coefficients respond to the sorting threshold \(\tau\) (near-fixation floor held
at \(\phi = 0.90\)). Because the per-unit sorting scores are stored, the
unidirectional set was recomputed at each \(\tau\) without re-running the sorting
test, and the standardisation was fixed over the full differentiated-unit set, so
the coefficients move only because \(\tau\) changes which units qualify.

\clearpage
\begin{center}
\includegraphics[width=\textwidth,trim=0 0 0 0.60in,clip]{../Figures/moduleA_direction_sweep.png}
\captionof{figure}{\textbf{Sorting direction is governed by the diagnostic index across sorting thresholds.}
Standardised coefficients (points; \(95\%\) confidence intervals) of the logistic
direction model P(\Faq)\,\(\sim\)\,diagnostic index \(+\) recombination \(+\)
parental MAF \(+\) \(\log\) block size, fitted to unidirectional LD-reduced units,
as a function of the sorting threshold \(\tau\) (near-fixation floor
\(\phi = 0.90\)). Panels, left to right: diagnostic index, recombination,
parental minor-allele frequency and \(\log\) block size; the dashed line marks a
null coefficient. The diagnostic index dominates at every threshold and
strengthens as \(\tau\) increases (from \(+1.57\) at \(\tau = 0.5\) to
\(+1.95\) at \(\tau = 0.8\)), whereas recombination contributes only a small,
stable negative coefficient (\(\approx -0.09\); low recombination leans \Faq{}
once the diagnostic index, parental MAF and block size are controlled) that
remains roughly an order of magnitude weaker than the diagnostic index
throughout, losing precision only at the strictest threshold as the
unidirectional set shrinks. Parental MAF and block size stay near zero.}
\label{figS:directionsweep}
\end{center}

% ---- Figure S4: genomic landscape of sorting across thresholds ---------
Finally, the sorting signal has clear genomic structure that persists across the
sorting threshold. Figure~\ref{figS:sortinglandscape} maps, along each chromosome,
the fraction of differentiated SNPs fixing toward each parent in 100-kb windows at
\(\tau \in \{0.5, 0.6, 0.7, 0.8\}\) (near-fixation floor \(\phi = 0.90\)), together with
the net direction, the diagnostic-index landscape and the marker-level local-LD
statistic \(\mathrm{ld}_w(0.95)\). Windows with fewer than 50 differentiated SNPs
are removed as low-density, since their per-window estimates are unstable.

\begin{landscape}
\begin{center}
\includegraphics[width=\linewidth]{../Figures/moduleA_sorting_manhattan_directional.png}
\captionof{figure}{\textbf{Genomic landscape of directional ancestry sorting across sorting thresholds.}
For each 100-kb window along the genome (faceted by chromosome), from top to
bottom: the fraction of differentiated SNPs fixing toward \Faq{}, the fraction
fixing toward \Fpo{}, the net direction (\Faq{} minus \Fpo{}), the diagnostic-index
landscape (grey points are individual markers, line is the 5-window rolling mean,
view clipped to \([-100, 0]\), dashed line marks the \(\mathrm{DI} = -25\)
diagnostic cut), and the marker-level local-LD statistic \(\mathrm{ld}_w(0.95)\)
(grey points, line is the window mean). Sorting rows are coloured by sorting
threshold \(\tau\) (near-fixation floor \(\phi = 0.90\); 5-window rolling mean).
Windows with fewer than 50 differentiated SNPs are removed as low-density. Sorting
peaks frequently overlapped regions of elevated local LD, whereas the diagnostic
index primarily predicted the direction of sorting. Both patterns persisted at more
stringent sorting thresholds.}
\label{figS:sortinglandscape}
\end{center}
\end{landscape}

\end{document}
